\section{Keeping Advent}

\subsection{An hour you can return to}

St Augustine writes to Proba about St Paul's command to pray without ceasing
(1 Thess. 5:17). Proba can spend an hour on business without ceasing to want
eternal life. When she keeps that desire before God, Augustine says, she
prays even while she is occupied. He also knows what business can do to
desire. We forget what we wanted most in the course of doing what had to be
done. So he bids her go to prayer at set times, when she can give words to
her desire and find it kindled again (Letter 130.18--20).

Keep a time you can find again tomorrow. For five minutes before work or
after supper, read the day's passage all the way through. Return to a
sentence that asks you to trust a promise, repent of a wrong, or thank God
for what he has given you. Bring that one thing to him. Say the Our Father
and remember someone by name. If you miss a day, return at your next usual opportunity and take up
that day's passage.

Some days are made by a child's sleep, a shift at work, or the needs of
someone ill. Fix the time by something that survives the changes. A parent
can open the psalm after clearing a place at the table. A worker can read
before entering the house at the end of a late shift. If you share the
practice with others, give someone else a turn to read, then leave room
for a little silence. A short prayer people can say together is enough.

For a fuller observance, spend fifteen minutes with the Scripture and that
week's reading from a Father in \emph{Readings for each day}. Take one line
away with you. Return to it around midday when you might have an occasion
to do what it asks. In the evening you can pray the Magnificat or the
\emph{Alma Redemptoris Mater} in \emph{Praying at home}. Before bed thank God
for a particular good you received and ask pardon for what you did wrong.

\subsection{When attention fails}

Psalm 24 gives a way to start. The psalmist asks, “Direct me in thy truth,
and teach me” (Ps. 24:5). Stay with the two requests for a little while.
What do I already know to be true that I would rather put out of mind?
Whose good counsel have I refused to hear? Then turn the psalmist's words
toward God. If something definite has come to you, such as the need to
speak honestly to somebody the next day, carry it away as a thing to do.

John Cassian records a conversation with Abbot Isaac in which one verse
serves from morning to night. When a monk's head droops over the book he
is trying to read, when old conversations interrupt prayer, and when
prayer comes easily and the monk begins to take credit for it, Isaac gives
him the same cry (\textit{Conferences} 10.10). “O God, come to my assistance;
O Lord, make haste to help me” (Ps. 69:2). You can say it as you sit down
to pray. Say it again when you discover you have been conducting an argument
in your head for the past three minutes. Begin where you left off.

Isaac chose a plea for God's help that can be spoken by the person
doing well and by the person beginning again. That is one reason to keep
a psalm close at hand all day. You can ask for the help to speak gently
before you open a door. When you come out of the room and know you spoke
badly, you can ask for the help to go back and put it right.

\subsection{The fast and the hungry}

Isaiah asks the worshipper to look at the people who bear the cost of
his life. A fast could be kept in sackcloth while men still oppressed
their workers. God asks them instead to break the burdens they have laid
on others and to share their bread. “Deal thy bread to the hungry,” the
prophet says, before promising that light will rise like morning (Isa.
58:5--8). The person who has laid the burden on someone else has a particular
work to do.

The fast also gives you something that can be handed over. On a day
each week when it is prudent, make a simpler meal and set aside what you
save. You could bring a bag of good food to the parish pantry or use that
money to help someone complete a week's shopping. If a change of diet is
imprudent, go without an optional purchase or give some of your free time
to someone who needs company or help. In his last words to Proba, Augustine
asks her household to bear with one another in charity while each prays,
fasts and keeps vigil according to what each can bear (Letter 130.31).

Keep the obligation of your own place and learn what its penitential
days ask. Advent fasting is voluntary under present Latin Church law.
Fridays retain their penitential character. The laws that governed when
the 1962 Missal appeared and the laws in force now, including United
States practice, are set out in \emph{About this book}.

\subsection{Receiving mercy}

Find a time for confession before the last crowded days. Look up the hours
or arrange a time with a priest. If you are conscious of grave sin, seek
confession promptly. Sit quietly beforehand and ask with the psalmist for
a clean heart (Ps. 50:12). Look at what you have done since your last
confession. Ask where you have left prayer or a duty to God aside, where
you have spoken against your neighbour or withheld what you owed, and
where you have let a habit make a choice for you. Remember the good you
have received as well. It helps to know whom you are asking for pardon.

In Jericho, Jesus stops under the tree where Zacchaeus has climbed to see
him. He tells the tax collector to come down because he means to stay in
his house. Zacchaeus makes haste and receives him with joy. In the house
he promises the poor half his goods and undertakes to pay back those he
has wronged. Jesus says salvation has come to that house (Luke 19:5--10).
The Lord asks to be received by a man whose dealings need to change.
He calls him out of the tree before the man has put his accounts in order.

When you confess, say plainly what needs forgiving. Ask for the grace to
turn from it. Through the priest's words of absolution Christ forgives,
exercising the mercy he entrusted to his apostles (John 20:22--23). If there is some wrong you can repair, think of the person
who has borne it. An overpayment can be returned. A story told at another
person's expense can be corrected with those who heard it. If you cannot
see what kind of repair would be just or possible, ask the confessor.
Do the penance he gives you. The Roman Catechism describes prayer,
fasting and alms as works of satisfaction that concern God, oneself
and one's neighbour (Part II, Penance). They are fitting ways for a
person who has received mercy to begin to live mercifully.
